% Generated draft: requires mathematical and source review.
\documentclass[11pt]{article}
\usepackage[margin=25mm]{geometry}
\usepackage{fontspec}
\setmainfont{Latin Modern Roman}
\usepackage{amsmath,amssymb,mathtools,unicode-math}
\setmathfont{Latin Modern Math}
\usepackage{enumitem,needspace,xurl,hyperref}
\setlength{\parindent}{0pt}
\setlength{\parskip}{5pt}
\setlength{\emergencystretch}{3em}
\newcommand{\N}{\mathbb{N}}
\newcommand{\Z}{\mathbb{Z}}
\newcommand{\Q}{\mathbb{Q}}
\newcommand{\R}{\mathbb{R}}
\newcommand{\C}{\mathbb{C}}
\newcommand{\F}{\mathbb{F}}
\newcommand{\A}{\mathbb{A}}
\newcommand{\PP}{\mathbb{P}}
\newcommand{\E}{\mathbb{E}}
\newcommand{\eps}{\varepsilon}
\newcommand{\dd}{\,\mathrm{d}}
\begin{document}
\section*{UnsolvedMath ID: 30003724}
\textbf{Source problem code:} catalogue entry
\subsection{Definitions and mathematical statement}
For an Abelian sandpile model in spatial dimension \(d\), once its underlying graph, configuration law, grain-addition rule, and stabilization procedure have been specified, let \(T_x\) denote the event that a vertex \(x\) topples at least once during that stabilization. Write
\[
p_d(x):=\mathbb P_d(T_x),
\]
where \(\mathbb P_d\) denotes the probability law of the specified experiment. This is conditional notation for a toppling probability; the supplied record does not specify the experiment or identify the toppling probability to which its question refers.

\textbf{Source question.} For the Abelian sandpile model, does the high-dimensional scaling limit for the toppling probability have an analogue in low dimensions \(d<5\)?

\textbf{Unresolved formulation.} The record does not state the high-dimensional limiting formula, the limiting parameter, the normalization, or the meaning of an analogue. It also does not specify the underlying graph, probability measure, addition site or rule, boundary conditions, or finite-volume versus infinite-volume formulation. Consequently, neither a particular limit involving \(p_d\) nor a precise quantified conjecture can be extracted without additional source information. The range \(d<5\) is retained exactly; the record does not explicitly enumerate the dimensions or state whether a single common formulation is intended for all of them. The background describes known high-dimensional toppling-probability asymptotics, while the question refers to a scaling limit; their intended relationship is unspecified.
\subsection{Short English statement}
The question asks whether the high-dimensional scaling behavior of sandpile toppling probabilities has a counterpart in dimensions \(d<5\). The supplied record leaves both the probability experiment and the proposed limiting statement unspecified, so a precise conjecture requires further source clarification.
\subsection{Sources}
\begin{itemize}
\item Local UnsolvedMath registry, supplied complete record for permanent ID \(30003724\), titled \emph{Low-Dimensional Sandpile Scaling Limits}; exact source problem code \texttt{catalogue entry}. The record belongs to the Oberwolfach Reports Open Problems collection and gives proposed year \(2018\); no individual proposer is supplied.
\item Supplied source citation: \emph{Strongly Correlated Random Interacting Processes} (2019), Oberwolfach Reports, DOI \texttt{10.4171/owr/2018/4}, pp.\ 187--253. Source link: \url{https://doi.org/10.4171/owr/2018/4}. The background describes Oberwolfach Reports as documenting mathematical workshops and their open-problem discussions.
\item Supplied literature source, labeled ``original OWR report'': \url{https://doi.org/10.4171/OWR/2018/4}.
\item Supplied literature source, labeled ``Hutchcroft, high-dimensional sandpile asymptotics'': \url{https://doi.org/10.1007/s00440-019-00923-3}.
\item Literature assessment dated exactly \texttt{2026-08-21}: the supplied curation labels the problem open and states that high-dimensional toppling-probability asymptotics are known, while an analogue in low dimensions was not established by the checked literature. This is supplied dated literature triage derived from the OWR curation, not an independently verified assessment; its time-sensitive status and cited claims require independent verification.
\item The registry assigns the statement status \texttt{corrected\_verified} and reports a manual semantic, typographic, and LaTeX review asserting that the clean statement is self-contained and preserves the target without unsupported additions. This is a supplied editorial qualification, not independent verification of the mathematical formulation or literature status.
\end{itemize}
\end{document}
